\section[Stage 002]{Stage 002: Breathing Reduction Back to the Old \((a,L)\) Closure}
\label{stage:002}

\stagefield{Part / anchor}{Part I; primary anchors \resultanchor{MTDC-T1}, \resultanchor{MTDC-T2}, and \resultanchor{MTDC-T4}.}
\claimstatus{\StatusExactClosure{} inside the Stage~001 quadratic wall closure.  Damping and open-system completion are not derived here.}

\stagefield{Purpose}{Stage~002 verifies that the distributed wall lift is a true lift of the old geometry sector rather than a replacement unrelated to it.  The axisymmetric lowest-mode truncation reproduces a two-coordinate \((\delta a,\delta L)\) matrix system, while the grouped real \(P_2\) modes appear as the next harmonic family of the same wall PDE.}

\stagefield{Inputs}{The inputs are the normalized monopole harmonic, the Stage~001 wall action, the axisymmetric split of \(\eta\), and the grouped real \(P_2\) harmonic eigenvalue \(-\Delta_{S^2}Y_{2m}=6Y_{2m}\).}

\stagefield{Verification}{SymPy audit: \StageFile{scripts/moving_throat_pde_stage002_breathing_reduction_sympy_audit.py}.  Mathematica audit: \StageFile{mathematica/moving_throat_pde_stage002_breathing_reduction_mathematica_audit.wl}.  The audits check the \(Y_{00}\) normalization bridge, the derived \(4\pi\) overlap factor, the conservative \((\delta a,\delta L)\) matrix reduction, and the uncoupled grouped real \(P_2\) degeneracy.}

\subsection*{Derivation}

\paragraph{1. Axisymmetric two-mode ansatz.}
Using
\begin{equation}
Y_{00}=\frac{1}{2\sqrt\pi},
\qquad
\int_{S^2}Y_{00}^2\,d\Omega=1,
\qquad
\frac{1}{4\pi}\int_{S^2}Y_{00}\,d\Omega=\frac{1}{2\sqrt\pi},
\end{equation}
the axisymmetric perturbation is truncated as
\begin{equation}
\label{eq:app-stage002-two-mode-ansatz}
\eta_0(w,t)=2\sqrt\pi\left[\alpha_a(w)\delta a(t)+\alpha_L(w)\delta L(t)\right].
\end{equation}
Set
\begin{equation}
Q^A=(\delta a,\delta L),
\qquad
A\in\{a,L\}.
\end{equation}

\paragraph{2. Reduction of the quadratic wall action.}
On the axisymmetric branch, the angular-stiffness term vanishes because \(l=0\).  In the densitized one-dimensional convention inherited from Stage~001, the background surface weight has already been absorbed into the effective axial coefficients and profiles, so the reduction is written with respect to \(dw\) rather than \(\sqrt{\gamma_0}\,dw\).  The action therefore reduces to a Lagrangian
\begin{equation}
\label{eq:app-stage002-lred}
L_{\rm red}^{(0)}
=\frac12 M_{AB}\dot Q^A\dot Q^B
-\frac12 K_{AB}Q^AQ^B,
\end{equation}
where
\begin{equation}
\label{eq:app-stage002-mass-matrix}
\boxed{
M_{AB}=4\pi\int dw\,\mu_\eta(w)\alpha_A(w)\alpha_B(w)},
\end{equation}
and
\begin{equation}
\label{eq:app-stage002-stiffness-matrix}
\boxed{
K_{AB}=4\pi\int dw\left[
T_w(w)\alpha_A'(w)\alpha_B'(w)+K_0(w)\alpha_A(w)\alpha_B(w)
\right]}.
\end{equation}
Here \(K_0(w)=K_\eta(w)\) on the scalar branch.

\paragraph{3. Euler--Lagrange equations.}
Varying \eqref{eq:app-stage002-lred} gives
\begin{equation}
\label{eq:app-stage002-el}
\boxed{
M_{AB}\ddot Q^B+K_{AB}Q^B=0}.
\end{equation}
This is the conservative linear version of the old two-coordinate geometry closure.  The old schematic form
\begin{equation}
M_{AB}\ddot Q^B+\Gamma_{AB}\dot Q^B=-\frac{\partial H_{\rm tot}}{\partial Q^A}
\end{equation}
is recovered only after adding open-system or effective damping data.  Stage~002 derives the conservative part, not the damping matrix.

\paragraph{4. Grouped real \(P_2\) one-mode reduction.}
For one real quadrupole component take
\begin{equation}
\label{eq:app-stage002-p2-ansatz}
\eta_{2m}(\Omega,w,t)=\beta_2(w)q_{2m}(t)Y_{2m}^{\rm real}(\Omega),
\qquad
-\Delta_{S^2}Y_{2m}^{\rm real}=6Y_{2m}^{\rm real}.
\end{equation}
Inserting into the same wall action gives
\begin{equation}
\label{eq:app-stage002-l2m}
L_{2m}=\frac12 M_2\dot q_{2m}^{\,2}-\frac12K_2q_{2m}^{\,2},
\end{equation}
with
\begin{equation}
\label{eq:app-stage002-m2}
\boxed{
M_2=\int dw\,\mu_\eta(w)\beta_2(w)^2},
\end{equation}
and
\begin{equation}
\label{eq:app-stage002-k2}
\boxed{
K_2=\int dw\left[
T_w(w)\beta_2'(w)^2+\bigl(K_\eta(w)+6T_\Omega(w)\bigr)\beta_2(w)^2
\right]}.
\end{equation}
Thus, before matter/gauge coupling or explicit anisotropy,
\begin{equation}
\label{eq:app-stage002-p2-eom}
\boxed{
M_2\ddot q_{2m}+K_2q_{2m}=0
}
\end{equation}
for every real \(P_2\) component.  This is the first degeneracy theorem for the grouped real quadrupole lane.

\stagefield{Output}{The stage outputs the \((\delta a,\delta L)\) matrix system \eqref{eq:app-stage002-el}, the effective matrices \eqref{eq:app-stage002-mass-matrix}--\eqref{eq:app-stage002-stiffness-matrix}, and the uncoupled grouped \(P_2\) mass/stiffness pair \eqref{eq:app-stage002-m2}--\eqref{eq:app-stage002-k2}.}

\stagefield{Checks}{\begin{verificationchecklist}
\item The factor \(4\pi\) in \eqref{eq:app-stage002-mass-matrix} and \eqref{eq:app-stage002-stiffness-matrix} follows from the \(2\sqrt\pi\) normalization in \eqref{eq:app-stage002-two-mode-ansatz}.
\item The \(l=2\) angular stiffness shift is \(l(l+1)T_\Omega=6T_\Omega\).
\item The grouped \(P_2\) degeneracy follows because \(K_2\) and \(M_2\) do not depend on the real component label on an isotropic reference throat.
\item The reduction preserves the old \((a,L)\) closure but does not derive the old damping matrix.
\end{verificationchecklist}}

\stagefield{Downstream use}{Stage~003 couples the Stage~002 wall modes to stable BdG support modes.  Stage~023 uses the same wall baseline as the \(K_A,M_A\) entries of the full grouped bundle.  Later geometry-contamination stages use the distinction between collective scalar geometry and grouped \(P_2\) wall modes.}
